\section{Reproducibility}
\label{app:repro}

Everything needed to rerun the study is in the \texttt{study/} directory of the repository:
\begin{itemize}
  \item \texttt{PREREGISTRATION.md}: hypotheses, metrics, decision rules and deviations;
  \item \texttt{corpus.toml}: repository URLs, release tags, commit SHAs, install and pytest arguments;
  \item \texttt{harness/}: cloning at the pinned commit, environment setup, raw-record writer;
  \item \texttt{experiments/}: one script per experiment (\texttt{e0\_baseline} to \texttt{e6\_flaky});
  \item \texttt{data/raw/}: append-only JSONL records (gzipped once a run has ended), one file per run, each carrying the project commit, runner commit, platform, Python version, environment hash (of \texttt{pip freeze}) and seed; every mutant result stores its exact patch;
  \item \texttt{analysis/}: scripts that produce \texttt{data/derived/}, the tables, the figures and \texttt{paper/generated/macros.tex}.
\end{itemize}

\noindent The full pipeline is:
\begin{verbatim}
cd study
make corpus      # clone and install the eight projects
make run         # E0 baseline, then H5, H1, H3, H4, H6, H2
make compress    # gzip the raw files of finished runs
make analyse     # derived CSVs, tables, figures, macros
make paper       # this PDF
\end{verbatim}

H5 and H2 measure time, so run them with nothing else running on the machine. The primary runs used an Apple-silicon Mac (macOS, arm64) with 10 cores and 16\,GB of memory, and 8 mutant workers.

\begin{table}[t]
\centering\small
\begin{tabular}{lrrrrrr}
\toprule
Project & H1 macOS & H1 Linux & H3 macOS & H3 Linux & H5 macOS & H5 Linux \\
\midrule
cachetools & 0.48 & 0.48 & 0.77 & 0.71 & 1.07 & 1.10 \\
toolz & 0.50 & 0.50 & 1.00 & 1.00 & 1.88 & 1.88 \\
\bottomrule
\end{tabular}
\caption{Linux replication on GitHub Actions (ubuntu-latest) against the macOS primary runs. H1: share of high-overlap pairs that are mutually non-dominated. H3: share of call-deletion survivors inside fully killed functions. H5: median per-block wall-time ratio, contexts over no coverage.}
\label{tab:replication}
\end{table}

